\section{Background and Motivation}
\label{sec:background}

% =====================================================================
% figs/cascade.tex — mechanism figure: the mispricing cascade under
% one-hop proxy costing. Parallel-panel design: SAME merge, two cost
% models; left grays out what the proxy cannot see.
% Requires TikZ libraries: positioning, arrows.meta (loaded in main.tex).
% =====================================================================
\begin{figure}[t]
\centering
\resizebox{\linewidth}{!}{%
\begin{tikzpicture}[
    font=\footnotesize,
    >={Stealth[length=1.5mm]},
    cell/.style={draw, rounded corners=1pt, minimum height=4.5mm,
                 minimum width=7mm, inner sep=2pt},
    mbff/.style={cell, thick, minimum height=8mm, minimum width=8mm},
    gate/.style={draw, minimum height=4mm, minimum width=4.5mm, inner sep=1pt},
    wire/.style={->, semithick},
    gcell/.style={cell, draw=gray!55, text=gray!55},
    ggate/.style={gate, draw=gray!55, text=gray!55},
    gwire/.style={->, semithick, gray!55},
    note/.style={inner sep=1pt, align=center}
  ]
  % ---------------- panel split ----------------
  \draw[dashed, gray!70] (4.35,0.55) -- (4.35,5.08);
  \node[note, anchor=west, font=\footnotesize\bfseries] at (0.0,4.90)
    {Proxy view (one hop)};
  \node[note, anchor=west, font=\footnotesize\bfseries] at (4.55,4.90)
    {Evaluator reality};

  % ============ LEFT: what the proxy prices ============
  \node[cell] (ff1) at (0.65,4.30) {FF$_1$};
  \node[cell] (ff2) at (0.65,3.14) {FF$_2$};
  \node[mbff] (mb)  at (2.45,3.72) {4b};
  \draw[wire] (ff1.east) -- ([yshift=1.7mm]mb.west);
  \draw[wire] (ff2.east) -- ([yshift=-1.7mm]mb.west);
  \node[note, gray] at (1.45,4.35) {merge};
  % the only wiring the proxy prices: the merged cell's local net
  \fill (3.60,3.72) circle (0.45mm);
  \draw (mb.east) -- (3.60,3.72);
  \node[note] at (3.30,4.30) {local net\\ (priced)};
  % downstream logic, grayed out: outside the one-hop model
  \node[ggate] (lg1)  at (2.45,2.34) {g$_1$};
  \node[gcell] (lffs) at (3.60,2.34) {FF$_s$};
  \draw[gwire] (mb.south) -- (lg1.north);
  \draw[gwire] (lg1.east) -- (lffs.west);
  \node[note, gray] at (1.75,1.62) {downstream logic:\\ not priced};
  \node[note, text=green!50!black] at (2.15,0.90)
    {est.\ $\Delta$cost $< 0$ $\rightarrow$ accept};

  % ============ RIGHT: what the evaluator charges ============
  \node[mbff, minimum width=7mm] (mb2) at (5.15,4.08) {4b};
  \node[gate] (g1)  at (6.35,4.08) {g$_1$};
  \node[gate] (g2)  at (7.25,4.08) {g$_2$};
  \node[cell] (ffs) at (8.30,4.08) {FF$_s$};
  \draw[wire] (mb2.east) -- (g1.west);
  \draw[wire] (g1.east)  -- (g2.west);
  \draw[wire] (g2.east)  -- (ffs.west);
  \node[note, text=red!70!black] at (7.05,3.40)
    {slack $< 0$ at FF$_s$ $\Rightarrow$ true $\Delta$cost $> 0$};
  % next decision inherits the corrupted state
  \node[cell, dashed, minimum width=13mm] (nm) at (5.55,2.00) {next merge};
  \draw[dashed, ->] (mb2.south) .. controls (5.15,2.95) .. (nm.north);
  \node[note, text width=26mm, anchor=west] at (6.70,1.92)
    {priced on the corrupted state; errors compound};
\end{tikzpicture}}
\caption{One merge, two prices: the one-hop proxy sees only the local net and accepts (left); the evaluator charges the downstream cone (right), and the next merge is priced on the corrupted state.}
\label{fig:cascade}
\end{figure}


\subsection{Problem~B and the Official Evaluator}
\label{sec:problemb}

The ICCAD 2024 CAD Contest Problem~B gives a legalized design: single- and multi-bit flip-flops, fixed combinational gates, and MBFF library variants differing in bit width, power, area, and pin geometry. A solution rebanks and replaces the flip-flops to minimize
\begin{equation}
S \;=\; \alpha \cdot \tns \;+\; \beta \cdot \mathit{Power} \;+\; \gamma \cdot \mathit{Area} \;+\; \lambda \cdot V,
\label{eq:score}
\end{equation}
where $V$ counts bins whose utilization exceeds a per-testcase threshold; legality (no overlaps, all cells on rows inside the die) is a hard constraint~\cite{yang2024contest,fang2024overview}. The weights $\alpha,\beta,\gamma,\lambda$ and a \emph{DisplacementDelay} constant vary per testcase, so the dominant term shifts between $\tns$ and power/area.

An official evaluator binary defines scoring, with timing semantics unlike conventional net models. Each arc's delay is the \emph{Displacement\-Delay} constant times the two-point half-perimeter wirelength from driver pin to individual sink pin, with no shared net bounding box or Steiner topology. Each sink's slack starts from its parse-time base value and moves by two deltas: the Q-pin delay change when the driving flip-flop changes cell type, and the wire-delay change of affected arcs. $\tns$ sums negative D-pin slacks over all flip-flops. A refinement engine seeking monotone score descent must reproduce these semantics, since Eq.~\eqref{eq:score} is evaluated on them alone.

Our flow builds on a top-ranked open-source implementation of the contest reference flow (Fig.~\ref{fig:framework}): debank MBFFs into single-bit registers, spread by slack-driven conjugate-gradient global placement, cluster by mean shift, bank, legalize onto rows, run detailed placement. We previously strengthened several stages: maximum-weight-matching pairing with proximity- and risk-weighted edge costs replaced greedy banking, and detailed placement gained cost-aware repairs. These enhancements keep the one-hop proxy pricing and are not claims of this paper; they make every refinement delta in Sec.~\ref{sec:experiments} conservative, since refinement is measured on the strengthened front end. Everything from the incremental oracle onward (Secs.~\ref{sec:oracle}--\ref{sec:operators}) is this paper's contribution.

\subsection{Compounding Mispricing under Proxy Cost Models}
\label{sec:proxy}

Refinement explores millions of candidate moves per benchmark, so per-candidate evaluator calls are infeasible; contest-style flows substitute an internal proxy. The baseline's proxy is representative: it reads stored slacks and re-prices one hop of wiring around the moved cell. Fig.~\ref{fig:cascade} shows why this cascades: the evaluator charges the merged cell through its downstream fanout, whose stored slacks the proxy never updates; each committed move thus corrupts the state that prices the next one. On the hardest benchmark, the compounded error leaves the one-hop model reporting final $\tns$ roughly 34\% below what the evaluator's score decomposition implies.

Two prior structural rebanking attempts in the same codebase measured the cascade. One matched flip-flops into higher-bit cells, one debanked then rebanked mispriced groups; they degraded the score by $+254\%$ and $+134\%$ on a power-dominated hidden case. Both generated the move class we later exploit profitably: merging smaller flip-flops into 4-bit cells. The moves were sound and the one-hop prices were wrong. These cascades motivate the prediction, confirmed in Sec.~\ref{sec:experiments}, that exact pricing turns the same move class into strict descent.

\subsection{Related Work}
\label{sec:related}

MBFF banking has a long post-placement lineage, from timing-constrained merging through flop-tray co-design to capacitated clustering at scale~\cite{chang2010postplacement,shyu2013effective,kahng2016floptray,kahng2024scalable}. Chen et al.~\cite{chen2024slackredist} redistribute positive slack along register chains to enlarge feasible clustering regions before banking. Despite the shared term, theirs is a pre-clustering slack-budget transformation; we re-price every move exactly after placement.

On Problem~B itself, the contest winner's method appears only as a two-page late-breaking result~\cite{li2025lbr} self-reporting a 0.988 average ratio against the first-place entry. The format precludes a full method description; to our knowledge no full-length paper has matched that level on this benchmark. Chen et al.~\cite{chen2026agile} adopt a banking-then-refinement flow close to ours but price moves with proxy models. TIMBER~\cite{dassarma2026timber} reports on the same testcases but is not evaluator-comparable (Sec.~\ref{sec:mainresults}). Gong et al.~\cite{gong2025iscas} target MBFF timing and power at advanced nodes without a public evaluator.

Timing-driven placers attack the same objective~\cite{liao2022dreamplace4,date2025tddp,romanvicharra2024ffcentric} but optimize proxy timing signals without closing the loop with the scoring function. Incremental STA engines~\cite{huang2015opentimer,huang2021opentimer2} maintain full-netlist timing under general delay models; our oracle is narrower and cheap enough for inner-loop pricing. Global placers handle density analytically~\cite{lu2015eplace,chen2008ntuplace3}. No published work priced the contest's bin-violation term $\lambda \cdot V$ in refinement together with exact timing; we price it per move alongside $\tns$.
